\section*{Capítulo \ref{ch:b2:integration} --- Integração}

\begin{solution}{exo:b2:integration:1}
$\int_0^1 \frac{\dd t}{\sqrt{t(1-t)}}$: perto de $0$, $\sim t^{-1/2}$
($\alpha = \frac12 < 1$: \omterm{def:b2:integration:improper}{converge}); perto de $1$, $\sim
(1-t)^{-1/2}$: \omterm{def:b2:integration:improper}{converge}. Convergente (seu valor é $\pi$, pela
substituição $t = \sin^2\theta$).

$\int_1^\infty \frac{\ln t}{t^2}$: $\frac{\ln t}{t^2} = o(t^{-3/2})$:
convergente (valor $1$ por partes).

$\int_0^\infty \frac{\dd t}{1 + t^2\sin^2 t}$: divergente. Perto de $t =
n\pi$, escreva $t = n\pi + u$: $\sin^2 t = \sin^2 u \leq u^2$, logo, em
$\abs u \leq \frac{1}{n}$, $1 + t^2\sin^2 t \leq 1 + (n\pi +
1)^2u^2 \leq C n^2 u^2 + 1$; portanto
\[
\int_{n\pi - 1/n}^{n\pi + 1/n} \frac{\dd t}{1 + t^2\sin^2 t}
\geq \int_{-1/n}^{1/n} \frac{\dd u}{1 + Cn^2u^2}
= \frac{2\arctan\sqrt C}{\sqrt C}\cdot\frac{1}{n} ,
\]
um termo de uma série divergente do tipo harmônico: somando sobre $n$, a
primitiva é ilimitada.
\end{solution}

\begin{solution}{exo:b2:integration:2}
Substitua $u = \lambda t$:
\[
\int_0^\infty t^n \eu^{-\lambda t}\dd t
= \frac{1}{\lambda^{n+1}}\int_0^\infty u^n\eu^{-u}\dd u
= \frac{\Gamma(n+1)}{\lambda^{n+1}} = \frac{n!}{\lambda^{n+1}} .
\]
Com $t = \eu^{-u}$ ($\dd t = -\eu^{-u}\dd u$):
\[
\int_0^1 (\ln t)^n \dd t = \int_0^{\infty} (-u)^n \eu^{-u}\,\dd u
= (-1)^n\, n! .
\]
\end{solution}

\begin{solution}{exo:b2:integration:3}
Convergência: o integrando é $\leq \frac{1}{1+t^2}$ perto de $\infty$
e limitado perto de $0$ (os dois fatores são minorados longe de $0$):
\omterm{def:b2:series:def}{absolutamente} convergente, para todo $x$. Substituindo $t = \frac1u$
($\dd t = -\frac{\dd u}{u^2}$):
\[
I(x) = \int_0^\infty \frac{1}{\bigl(1 +
\frac1{u^2}\bigr)\bigl(1 + u^{-x}\bigr)}\cdot\frac{\dd u}{u^2}
= \int_0^\infty \frac{u^x}{(1 + u^2)(1 + u^x)}\,\dd u .
\]
Somando as duas expressões de $I(x)$:
\[
2I(x) = \int_0^\infty \frac{1 + t^x}{(1+t^2)(1+t^x)}\dd t
= \int_0^\infty \frac{\dd t}{1 + t^2} = \frac{\pi}{2} :
\]
$I(x) = \frac\pi4$, independente de $x$.
\end{solution}

\begin{solution}{exo:b2:integration:4}
Perto de $0^+$, com $u = \abs{\ln t} \to \infty$. Se $\alpha < 1$:
convergência qualquer que seja $\beta$ (compare com $t^{-\alpha'}$ para
$\alpha < \alpha' < 1$: o fator logarítmico é vencido). Se $\alpha > 1$:
divergência qualquer que seja $\beta$ (compare com $t^{-\alpha''}$,
$1 < \alpha'' < \alpha$). Se $\alpha = 1$: substitua $t =
\eu^{-u}$:
\[
\int_0^{1/2} \frac{\dd t}{t\,\abs{\ln t}^\beta}
= \int_{\ln 2}^{\infty} \frac{\dd u}{u^\beta},
\]
convergente se e somente se $\beta > 1$. Resumo: convergência se e somente se $\alpha < 1$,
ou ($\alpha = 1$ e $\beta > 1$) --- o espelho das séries
de Bertrand.
\end{solution}

\begin{solution}{exo:b2:integration:5}
\omterm{def:b2:metric:continuity}{Continuidade} em $\intco{0}{\infty}$: dominação
$\bigl|\frac{\eu^{-xt}}{1+t^2}\bigr| \leq \frac{1}{1+t^2}$,
integrável, uniforme em $x \geq 0$:
\cref{thm:b2:integration:continuity}.

$C^2$ em $\intoo{0}{\infty}$: em $x \geq a > 0$, as duas primeiras
derivadas em $x$, $\frac{-t\,\eu^{-xt}}{1+t^2}$ e
$\frac{t^2\eu^{-xt}}{1+t^2}$, são dominadas por $t\,\eu^{-at}$ e
$\eu^{-at}$: duas aplicações do
\cref{thm:b2:integration:leibnizrule}. Então
\[
F''(x) + F(x) = \int_0^\infty \frac{t^2 + 1}{1 + t^2}\,\eu^{-xt}\dd
t = \int_0^\infty \eu^{-xt}\dd t = \frac1x .
\]
Limite: $0 \leq F(x) \leq \int_0^\infty \eu^{-xt}\dd t = \frac1x \to
0$.
\end{solution}

\begin{solution}{exo:b2:integration:6}
Em $\intcc{\varepsilon}{M}$, substitua $u = at$ e $u = bt$ nas
duas metades:
\[
\int_\varepsilon^M \frac{f(at) - f(bt)}{t}\dd t
= \int_{a\varepsilon}^{aM}\frac{f(u)}{u}\dd u
- \int_{b\varepsilon}^{bM}\frac{f(u)}{u}\dd u
= \int_{a\varepsilon}^{b\varepsilon} \frac{f(u)}{u}\dd u
- \int_{aM}^{bM} \frac{f(u)}{u}\dd u .
\]
Primeira peça: $f(u) = f(0) + o(1)$ perto de $0$, e $\int_{a\varepsilon}
^{b\varepsilon} \frac{\dd u}{u} = \ln\frac ba$: a peça tende a
$f(0)\ln\frac ba$. Segunda peça: $f(u) \to f(\infty)$, mesmo
cálculo: tende a $f(\infty)\ln\frac ba$. Logo a \omterm{def:b2:integration:improper}{integral imprópria}
\omterm{def:b2:integration:improper}{converge} para $\bigl(f(0) - f(\infty)\bigr)\ln\frac ba$.

Com $f(t) = \eu^{-t}$ ($f(0) = 1$, $f(\infty) = 0$), $a = 1$, $b
= 2$:
\[
\int_0^\infty \frac{\eu^{-t} - \eu^{-2t}}{t}\dd t = \ln 2 .
\]
\end{solution}

\begin{solution}{exo:b2:integration:7}
Estenda o integrando por $0$ além de $t = n$: $g_n(t) = (1 -
\frac tn)^n t^{x-1}\mathbf{1}_{t \leq n}$. Pontualmente, $g_n(t) \to
\eu^{-t}t^{x-1}$ (o limite dos juros compostos, volume do primeiro ano de graduação).
Dominação: $\ln(1 - u) \leq -u$ dá $(1 - \frac tn)^n \leq
\eu^{-t}$ em $\intcc{0}{n}$, logo $\abs{g_n(t)} \leq
\eu^{-t}t^{x-1} = \varphi(t)$, integrável. Convergência dominada:
\[
\int_0^n \Bigl(1 - \frac tn\Bigr)^n t^{x-1}\dd t
\xrightarrow[n\to\infty]{} \int_0^\infty \eu^{-t}t^{x-1}\dd t
= \Gamma(x) .
\]
(Calcular o membro da esquerda por partes repetidas dá a forma de produto
de Euler $\Gamma(x) = \lim \frac{n!\,n^x}{x(x+1)\cdots(x+n)}$.)
\end{solution}

\begin{solution}{exo:b2:integration:8}
$H$ é derivável em $x$ (integrando $C^1$ em $x$, derivada
$-2x(1+t^2)\cdot\frac{\eu^{-x^2(1+t^2)}}{1+t^2} =
-2x\,\eu^{-x^2}\eu^{-x^2t^2}$, \omterm{def:b2:metric:continuity}{contínua} e limitada nos \omterm{def:b2:metric:compact}{compactos}
de $x$, dominação sobre $t \in \intcc{0}{1}$ trivial):
\[
H'(x) = -2x\,\eu^{-x^2}\int_0^1 \eu^{-x^2t^2}\,\dd t
\overset{u = xt}{=} -2\,\eu^{-x^2}\int_0^x \eu^{-u^2}\,\dd u
= -G'(x),
\]
pois $G'(x) = 2\eu^{-x^2}\int_0^x \eu^{-t^2}\dd t$ (regra da cadeia no
quadrado, teorema fundamental do cálculo). Logo $G + H$ é
constante, igual a $G(0) + H(0) = 0 + \int_0^1 \frac{\dd t}{1+t^2}
= \frac\pi4$.

Quando $x \to \infty$: $0 \leq H(x) \leq \eu^{-x^2}\int_0^1 \dd t \to
0$, logo $G(x) \to \frac\pi4$:
\[
\int_0^\infty \eu^{-t^2}\dd t = \sqrt{\frac\pi4} =
\frac{\sqrt\pi}{2} .
\]
(Por consequência, $\Gamma\bigl(\frac12\bigr) = 2\int_0^\infty
\eu^{-t^2}\dd t = \sqrt\pi$, pela substituição $t = \sqrt u$.)
\end{solution}

\begin{solution}{exo:b2:integration:9}
Cauchy--Schwarz (volume do primeiro ano de graduação, válida em $\intcc{\varepsilon}{M}$
e passada ao limite) aplicada à fatoração
$t^{\frac{x+y}{2}-1}\eu^{-t} = \bigl(t^{x-1}\eu^{-t}\bigr)^{1/2}
\bigl(t^{y-1}\eu^{-t}\bigr)^{1/2}$:
\[
\Gamma\Bigl(\frac{x+y}{2}\Bigr) \leq
\Gamma(x)^{1/2}\,\Gamma(y)^{1/2}
\quad\Longrightarrow\quad
\ln\Gamma\Bigl(\frac{x+y}{2}\Bigr) \leq \frac{\ln\Gamma(x) +
\ln\Gamma(y)}{2} :
\]
$\ln\Gamma$ é convexa no ponto médio; sendo \omterm{def:b2:metric:continuity}{contínua}
(\cref{thm:b2:integration:gammaprops}), ela é convexa
(\cref{exo:b2:realfun:8}). (A log-convexidade fixa $\Gamma$
de modo único entre as interpolações do fatorial --- o
teorema de Bohr--Mollerup, uma pérola do terceiro ano.)
\end{solution}

\begin{solution}{exo:b2:integration:10}
\begin{enumerate}
  \item Em $x \geq a > 0$: a derivada em $x$ do integrando é
        $-\sin t\,\eu^{-xt}$, dominada por $\eu^{-at}$:
        o~\cref{thm:b2:integration:leibnizrule} dá $F'(x) =
        -\int_0^\infty \eu^{-xt}\sin t\,\dd t$. Duas integrações por
        partes (ou a exponencial complexa):
        \[
        \int_0^\infty \eu^{-xt}\sin t\,\dd t
        = \Im \int_0^\infty \eu^{(-x+\iu)t}\dd t
        = \Im\frac{1}{x - \iu} = \frac{1}{1 + x^2} .
        \]
  \item $\abs{F(x)} \leq \int_0^\infty \eu^{-xt}\dd t = \frac1x \to
        0$. Integrando $F' = -\frac{1}{1+x^2}$ de $x$ a
        $\infty$: $0 - F(x) = -\bigl(\frac\pi2 - \arctan x\bigr)$,
        logo $F(x) = \frac\pi2 - \arctan x$.
  \item Fazendo $x \to 0^+$ com a \omterm{def:b2:metric:continuity}{continuidade} admitida: $F(0^+) =
        \frac\pi2$, e $F(0) = \int_0^\infty \frac{\sin t}{t}\dd
        t$ (a integral semiconvergente de Dirichlet,
        \cref{ex:b2:integration:sint}): seu valor é
        $\frac\pi2$.
\end{enumerate}
\end{solution}

\begin{solution}{exo:b2:integration:11}
Convergência: perto de $0$ o integrando se estende \omterm{def:b2:metric:continuity}{continuamente} pelo
valor $1$ ($\sin t \sim t$); no infinito ele é $\leq t^{-2}$:
convergência absoluta. Em $\intcc{\varepsilon}{M}$, integre por
partes com $u = \sin^2 t$, $v' = t^{-2}$:
\[
\int_\varepsilon^M \frac{\sin^2 t}{t^2}\dd t
= \Bigl[-\frac{\sin^2 t}{t}\Bigr]_\varepsilon^M
+ \int_\varepsilon^M \frac{2\sin t\cos t}{t}\dd t
= \Bigl[-\frac{\sin^2 t}{t}\Bigr]_\varepsilon^M
+ \int_{2\varepsilon}^{2M} \frac{\sin u}{u}\dd u
\]
($u = 2t$ na última integral). O colchete tende a $0$ nas duas
pontas ($\sin^2\varepsilon/\varepsilon \leq \varepsilon$;
$\sin^2 M/M \leq 1/M$), e a última integral tende a
$\int_0^\infty \frac{\sin u}{u}\dd u = \frac\pi2$
(\cref{exo:b2:integration:10}). Logo
\[
\int_0^\infty \Bigl(\frac{\sin t}{t}\Bigr)^{\!2}\dd t
= \frac\pi2 .
\]
\end{solution}

\begin{solution}{exo:b2:integration:12}
\begin{enumerate}
  \item Partes com $u = \frac{-1}{2t}$, $v' = -2t\,\eu^{-t^2}$
        (de modo que $v = \eu^{-t^2}$):
        \[
        T(x) = \Bigl[\frac{-\eu^{-t^2}}{2t}\Bigr]_x^\infty
        - \int_x^\infty \frac{\eu^{-t^2}}{2t^2}\dd t
        = \frac{\eu^{-x^2}}{2x}
        - \int_x^\infty \frac{\eu^{-t^2}}{2t^2}\dd t .
        \]
        Mesmo dispositivo na nova integral ($u = \frac{-1}{4t^3}$,
        $v' = -2t\,\eu^{-t^2}$):
        \[
        \int_x^\infty \frac{\eu^{-t^2}}{2t^2}\dd t
        = \frac{\eu^{-x^2}}{4x^3}
        - \frac34\int_x^\infty \frac{\eu^{-t^2}}{t^4}\dd t ,
        \]
        donde a identidade anunciada.
  \item Mais uma integração por partes majora o resto:
        \[
        \int_x^\infty \frac{\eu^{-t^2}}{t^4}\dd t
        = \frac{\eu^{-x^2}}{2x^5}
        - \frac52\int_x^\infty\frac{\eu^{-t^2}}{t^6}\dd t
        \leq \frac{\eu^{-x^2}}{2x^5},
        \]
        logo $0 \leq \frac34\int_x^\infty t^{-4}\eu^{-t^2}\dd t
        \leq \frac{3}{8x^5}\eu^{-x^2}$. Descartar o resto (positivo)
        na identidade da questão 1 dá a cota
        inferior; descartar o segundo termo (negativo) da primeira
        integração por partes dá $T(x) \leq \frac{\eu^{-x^2}}{2x}$. Dividindo
        o enquadramento por $\frac{\eu^{-x^2}}{2x}$: a razão fica
        espremida entre $1 - \frac{1}{2x^2}$ e $1$, logo $T(x)
        \sim \frac{\eu^{-x^2}}{2x}$.
  \item Iterar as partes produz a série formal
        \[
        T(x) \approx \frac{\eu^{-x^2}}{2x}\Bigl(1 -
        \frac{1}{2x^2} + \frac{1\cdot3}{(2x^2)^2} -
        \frac{1\cdot3\cdot5}{(2x^2)^3} + \cdots\Bigr),
        \]
        cujo $k$-ésimo coeficiente $1\cdot3\cdots(2k-1) =
        \frac{(2k)!}{2^k k!}$ cresce mais depressa do que qualquer sequência
        geométrica: para $x$ fixo os termos
        $\frac{1\cdot3\cdots(2k-1)}{(2x^2)^k}$ tendem a infinito
        (sua razão é $\frac{2k+1}{2x^2} \to \infty$), de modo que a
        série diverge para todo $x$. Trata-se de um
        desenvolvimento \emph{assintótico}: truncado em qualquer \omterm{def:b2:structures:generated}{ordem}
        fixa, o erro é da ordem do primeiro termo omitido
        quando $x \to \infty$ --- mas nunca de uma série
        convergente. (Essa estimativa de cauda é a cota padrão da cauda
        gaussiana dos capítulos de probabilidade.)
\end{enumerate}
\end{solution}

\begin{solution}{pb:b2:integration:1}
\textbf{1.} Em $0^+$ o integrando é $\sim t^{x-1}$: a
escala de extremidade finita \omterm{def:b2:integration:improper}{converge} se e somente se $1 - x < 1$, isto é, $x > 0$
(e, para $x \leq 0$, $t^{x-1} \geq t^{-1}$ diverge); em
$+\infty$, $t^{x-1}\eu^{-t} = o(t^{-2})$ \omterm{def:b2:integration:improper}{converge} para todo
$x$. Então $\Gamma(x) = \frac{\Gamma(x+1)}{x}$ e $\Gamma(x+1)
\to \Gamma(1) = 1$ quando $x \to 0^+$ (\omterm{def:b2:metric:continuity}{continuidade},
\cref{thm:b2:integration:gammaprops}): $\Gamma(x) \sim
\frac1x$.

\textbf{2.} Com $t = u^2$, $\dd t = 2u\,\dd u$:
\[
\Gamma\Bigl(\frac12\Bigr) = \int_0^\infty
t^{-1/2}\eu^{-t}\dd t
= \int_0^\infty \frac{\eu^{-u^2}}{u}\,2u\,\dd u
= 2\int_0^\infty \eu^{-u^2}\dd u = \sqrt\pi
\]
pelo~\cref{exo:b2:integration:8}. Com $u = v/\sqrt2$:
\[
\int_\R \eu^{-v^2/2}\dd v
= 2\sqrt2\int_0^\infty \eu^{-u^2}\dd u
= \sqrt2\,\sqrt\pi = \sqrt{2\pi} .
\]

\textbf{3.} Verdadeiro para $n = 0$ (os dois lados valem $\sqrt\pi$). Se
$\Gamma(n + \frac12) = \frac{(2n)!}{4^n n!}\sqrt\pi$, a
equação funcional dá
\[
\Gamma\Bigl(n + 1 + \frac12\Bigr)
= \Bigl(n + \frac12\Bigr)\Gamma\Bigl(n + \frac12\Bigr)
= \frac{2n+1}{2}\cdot\frac{(2n)!}{4^n n!}\sqrt\pi
= \frac{(2n+2)!}{4^{n+1}(n+1)!}\sqrt\pi ,
\]
o último passo porque $\frac{(2n+2)!}{(2n)!} = (2n+2)(2n+1)$
e $\frac{2n+1}{2} = \frac{(2n+2)(2n+1)}{4(n+1)}$.

\textbf{4.} O~\cref{thm:b2:integration:gammaprops} dá
$\Gamma''(x) = \int_0^\infty t^{x-1}\eu^{-t}(\ln t)^2\dd t$
(duas aplicações da regra de Leibniz, com dominações como na
demonstração do teorema); o integrando é $\geq 0$ e não identicamente
nulo, logo $\Gamma'' > 0$: $\Gamma$ é estritamente convexa e
$\Gamma'$ é estritamente crescente. Como $\Gamma(1) = \Gamma(2)
= 1$, Rolle fornece $x_0 \in \intoo12$ com $\Gamma'(x_0) =
0$; a monotonicidade estrita de $\Gamma'$ faz de $x_0$ seu único
zero, com $\Gamma' < 0$ antes e $\Gamma' > 0$ depois:
$\Gamma$ decresce em $\intoo0{x_0}$, cresce em
$\intoo{x_0}\infty$, e $x_0$ é o mínimo único.

\textbf{5.} Sejam $k \in \N$ e $x \geq 3$; escolha o inteiro
$n$ com $n + 1 \leq x < n + 2$ (de modo que $n \geq 1$). Pela
monotonicidade da questão 4 (válida a partir de $x_0 < 2$): $\Gamma(x)
\geq \Gamma(n + 1) = n!$, enquanto $x^k \leq (n+2)^k$. Como
$\frac{n!}{(n+2)^k} \to \infty$ (os fatoriais vencem as potências, volume do primeiro ano
de graduação), $\frac{\Gamma(x)}{x^k} \geq \frac{n!}{(n+2)^k} \to
\infty$ quando $x \to \infty$: $x^k = o(\Gamma(x))$.

\textbf{6.} Perto de $0$ o integrando é $\sim t^{x-1}$
(convergente se e somente se $x > 0$), perto de $1$ ele é $\sim (1-t)^{y-1}$
(se e somente se $y > 0$); as duas comparações são entre funções positivas,
logo $B(x,y)$ \omterm{def:b2:integration:improper}{converge} exatamente para $x, y > 0$. A substituição
$t \mapsto 1 - t$ troca os dois fatores: $B(x,y) = B(y,x)$.

\textbf{7.} $B(x,1) = \int_0^1 t^{x-1}\dd t = \frac1x$. Partes
em $\intcc\varepsilon{1-\varepsilon}$ com $u = (1-t)^y$, $v =
\frac{t^x}{x}$:
\[
\int t^{x-1}(1-t)^{y}\dd t
= \Bigl[\frac{t^x(1-t)^y}{x}\Bigr]
+ \frac{y}{x}\int t^{x}(1-t)^{y-1}\dd t ;
\]
o colchete se anula nas duas pontas quando $\varepsilon \to 0$ ($x >
0$ em $0$, $y > 0$ em $1$), restando $B(x, y+1) = \frac
yx\,B(x+1, y)$.

\textbf{8.} Como $t + (1-t) = 1$:
\[
t^{x-1}(1-t)^{y-1} = t^{x}(1-t)^{y-1} + t^{x-1}(1-t)^{y},
\]
logo $B(x,y) = B(x+1,y) + B(x,y+1)$. A questão 7 se lê $B(x+1,y) =
\frac xy B(x,y+1)$; substituindo,
\[
B(x,y) = \Bigl(\frac xy + 1\Bigr)B(x,y+1)
= \frac{x+y}{y}\,B(x,y+1),
\]
isto é, $B(x,y+1) = \frac{y}{x+y}B(x,y)$; a relação gêmea
segue da simetria da questão 6.

\textbf{9.} Indução sobre $n$ com $m$ fixo: $B(m,1) = \frac1m =
\frac{(m-1)!\,0!}{m!}$ e, se a fórmula vale em $n$,
\[
B(m, n+1) = \frac{n}{m+n}\,B(m,n)
= \frac{n}{m+n}\cdot\frac{(m-1)!(n-1)!}{(m+n-1)!}
= \frac{(m-1)!\,n!}{(m+n)!} .
\]
Reescrevendo: $B(m,n) = \frac{(m-1)!(n-1)!}{(m+n-1)!} =
\bigl[(m+n-1)\binom{m+n-2}{m-1}\bigr]^{-1}$.

\textbf{10.} Os dois lados de $B(x,n) =
\frac{\Gamma(x)\Gamma(n)}{\Gamma(x+n)}$ valem $\frac1x$ em $n =
1$ ($\Gamma(1) = 1$, $\Gamma(x+1) = x\Gamma(x)$). Se eles coincidem
em $n$, então, pela relação de descida e pela equação
funcional:
\[
B(x, n+1) = \frac{n}{x+n}\,B(x,n), \qquad
\frac{\Gamma(x)\Gamma(n+1)}{\Gamma(x+n+1)}
= \frac{n}{x+n}\cdot
\frac{\Gamma(x)\Gamma(n)}{\Gamma(x+n)} :
\]
as duas sequências obedecem à mesma recursão a partir da mesma semente,
logo coincidem para todo $n \in \N^*$ e todo $x > 0$.

\textbf{11.} Com $t = \sin^2\theta$ ($\theta \in
\intoo0{\pi/2}$, $\dd t = 2\sin\theta\cos\theta\,\dd\theta$),
$t^{x-1} = \sin^{2x-2}\theta$ e $(1-t)^{y-1} =
\cos^{2y-2}\theta$:
\[
B(x,y) = \int_0^{\pi/2}\sin^{2x-2}\theta\,\cos^{2y-2}\theta
\cdot 2\sin\theta\cos\theta\,\dd\theta
= 2\int_0^{\pi/2}\sin^{2x-1}\theta\,\cos^{2y-1}\theta\,
\dd\theta .
\]

\textbf{12.} Tome $y = \frac12$ (matando o fator cosseno) e
$2x - 1 = n$: $B\bigl(\frac{n+1}2, \frac12\bigr) = 2W_n$, isto é,
$W_n = \frac12 B\bigl(\frac{n+1}2,\frac12\bigr)$. A relação de descida
na primeira variável dá
\[
\frac{W_n}{W_{n-2}}
= \frac{B\bigl(\frac{n-1}2 + 1, \frac12\bigr)}
{B\bigl(\frac{n-1}2, \frac12\bigr)}
= \frac{\frac{n-1}2}{\frac{n-1}2 + \frac12}
= \frac{n-1}{n} :
\]
a recorrência de Wallis, desta vez sem nenhuma integração por partes
sobre senos --- a Parte II fez o trabalho de uma vez por todas.

\textbf{13.} $B\bigl(\frac12,\frac12\bigr) = 2W_0 =
2\cdot\frac\pi2 = \pi$, enquanto
$\Gamma\bigl(\frac12\bigr)^2/\Gamma(1) = (\sqrt\pi)^2 = \pi$:
a fórmula de Euler vale em $\bigl(\frac12,\frac12\bigr)$.

\textbf{14.} Iterando $W_{2n} = \frac{2n-1}{2n}W_{2n-2}$ a partir de
$W_0 = \frac\pi2$:
\[
W_{2n} = \frac\pi2\prod_{k=1}^{n}\frac{2k-1}{2k}
= \frac\pi2\cdot\frac{(2n)!}{4^n(n!)^2},
\]
pois $\prod(2k-1) = \frac{(2n)!}{2^n n!}$ e $\prod 2k = 2^n
n!$. Logo, usando a questão 3:
\[
B\Bigl(n+\frac12, \frac12\Bigr) = 2W_{2n}
= \pi\,\frac{(2n)!}{4^n(n!)^2}
= \frac{(2n)!\sqrt\pi}{4^n n!}\cdot\frac{\sqrt\pi}{n!}
= \frac{\Gamma\bigl(n+\frac12\bigr)\Gamma\bigl(\frac12\bigr)}
{\Gamma(n+1)} .
\]
Fixe agora $x \in \frac12\N^*$. A fórmula de Euler vale em $(x,
\frac12)$: para $x$ inteiro isso é a questão 10 (com a simetria),
para $x = n + \frac12$ é a fórmula acima. Os dois lados da
fórmula de Euler obedecem à recursão de descida $y \mapsto y + 1$
(questão 8 à esquerda, equação funcional à direita,
como na questão 10): a indução propaga a fórmula de $y =
\frac12$ e $y = 1$ a todo $y \in \frac12\N^*$. A fórmula de
Euler vale, portanto, sempre que $2x, 2y \in \N^*$.

\textbf{15.} Com $u = \frac{t}{1-t}$, isto é, $t =
\frac{u}{1+u}$, $1 - t = \frac{1}{1+u}$, $\dd t =
\frac{\dd u}{(1+u)^2}$:
\[
B(x,y) = \int_0^\infty
\Bigl(\frac{u}{1+u}\Bigr)^{x-1}
\Bigl(\frac{1}{1+u}\Bigr)^{y-1}
\frac{\dd u}{(1+u)^2}
= \int_0^\infty \frac{u^{x-1}}{(1+u)^{x+y}}\,\dd u .
\]
Em $x = y = \frac12$, com $u = v^2$:
\[
\int_0^\infty \frac{u^{-1/2}}{1+u}\dd u
= \int_0^\infty \frac{2\,\dd v}{1+v^2} = \pi
= B\Bigl(\frac12,\frac12\Bigr) . \checkmark
\]

\textbf{16.} Uma integração por partes, para $1 \leq k \leq n$
e $s > 0$ ($u = (1 - t/n)^k$, $v = t^s/s$; os termos de bordo
se anulam):
\[
\int_0^n \Bigl(1-\frac tn\Bigr)^{\!k} t^{s-1}\dd t
= \frac{k}{ns}\int_0^n \Bigl(1-\frac tn\Bigr)^{\!k-1}
t^{s}\dd t .
\]
Partindo de $k = n$, $s = x$ e iterando $n$ vezes:
\[
\int_0^n \Bigl(1-\frac tn\Bigr)^{\!n} t^{x-1}\dd t
= \frac{n(n-1)\cdots1}{n^n\,x(x+1)\cdots(x+n-1)}
\int_0^n t^{x+n-1}\dd t
= \frac{n!}{n^n}\cdot
\frac{n^{x+n}}{x(x+1)\cdots(x+n)} ,
\]
que é $\dfrac{n!\,n^x}{x(x+1)\cdots(x+n)}$.

\textbf{17.} Pelo~\cref{exo:b2:integration:7} o membro da esquerda tende
a $\Gamma(x)$ (convergência dominada com dominante
$t^{x-1}\eu^{-t}$); o membro da direita é o quociente de Gauss:
\[
\Gamma(x) = \lim_{n\to\infty}
\frac{n!\,n^x}{x(x+1)\cdots(x+n)} .
\]

\textbf{18.} Tomando logaritmos no quociente da questão 16,
$G_n(x)$, e separando $\ln(x+k) = \ln k + \ln(1 + x/k)$ para $k
\geq 1$:
\[
\ln G_n(x) = x\ln n - \ln x - \sum_{k=1}^n
\ln\Bigl(1+\frac xk\Bigr)
= -\ln x + x(\ln n - H_n)
+ \sum_{k=1}^n\Bigl(\frac xk -
\ln\Bigl(1+\frac xk\Bigr)\Bigr).
\]
Para $u \geq 0$, $u - \frac{u^2}2 \leq \ln(1+u) \leq u$, de modo que o
termo geral está em $\intcc{0}{x^2/(2k^2)}$: a série
\omterm{def:b2:integration:improper}{converge} (comparação com $\sum k^{-2}$). Como $\ln n - H_n
\to -\gamma$ (\cref{ex:b2:comparison:harmonic}) e $\ln G_n(x)
\to \ln\Gamma(x)$ (questão 17 e \omterm{def:b2:metric:continuity}{continuidade} de $\ln$):
\[
\ln\Gamma(x) = -\ln x - \gamma x
+ \sum_{k=1}^\infty\Bigl(\frac xk -
\ln\Bigl(1+\frac xk\Bigr)\Bigr) .
\]

\textbf{19.} Em $x = \frac12$, o denominador é
$\prod_{k=0}^n\bigl(k+\frac12\bigr) =
\frac{(2n+1)!}{2^{2n+1}n!}$ (desenvolvendo as metades), logo
\[
G_n\Bigl(\frac12\Bigr)
= \frac{n!\,\sqrt n\;2^{2n+1}n!}{(2n+1)!}
= \frac{2\sqrt n\;4^n}{(2n+1)\binom{2n}{n}} .
\]
Com $\binom{2n}{n} \sim \frac{4^n}{\sqrt{\pi n}}$
(\cref{ex:b2:comparison:centralbinomial}):
\[
G_n\Bigl(\frac12\Bigr)
\sim \frac{2\sqrt n\,\sqrt{\pi n}}{2n+1}
\longrightarrow \sqrt\pi = \Gamma\Bigl(\frac12\Bigr) .
\]
O $\sqrt\pi$ do coeficiente binomial central (que veio de
Wallis, logo da constante de Stirling) e o $\sqrt\pi$ da integral
gaussiana são o mesmo número.

\textbf{20.} A identidade é álgebra direta: multiplique
$\frac{n!\,n^x}{x(x+1)\cdots(x+n)}$ por $\frac{nx}{x+n+1}$ e
absorva $x$ no produto e $n$ em $n^x$. Fazendo $n \to
\infty$: o membro da esquerda tende a $\Gamma(x+1)$ (Gauss em $x+1$),
o da direita a $\Gamma(x)\cdot x\cdot 1$, pois
$\frac{n}{x+n+1} \to 1$: $\Gamma(x+1) = x\Gamma(x)$ ---
recuperado sem uma única integração por partes.

\textbf{21.} Com $u = t^a$, $t = u^{1/a}$, $\dd t = \frac1a
u^{1/a - 1}\dd u$:
\[
\int_0^\infty \eu^{-t^a}\dd t
= \frac1a\int_0^\infty u^{\frac1a - 1}\eu^{-u}\dd u
= \frac1a\,\Gamma\Bigl(\frac1a\Bigr)
= \Gamma\Bigl(1 + \frac1a\Bigr) .
\]
Quando $a \to +\infty$ (ao longo de qualquer sequência): $\eu^{-t^a} \to 1$ para
$0 < t < 1$, $\to \eu^{-1}$ em $t = 1$, $\to 0$ para $t > 1$;
para $a \geq 2$ domine por $\mathbf 1_{t \leq 1} +
\eu^{-t^2}\mathbf 1_{t > 1}$ ($t^a \geq t^2$ para $t \geq 1$),
integrável. Convergência dominada: a integral tende a
$\int_0^1 1\,\dd t = 1$ --- como tinha de ser, pois $\Gamma(1 +
\frac1a) \to \Gamma(1) = 1$ pela \omterm{def:b2:metric:continuity}{continuidade}.

\textbf{22.} Com $u = t^n$, $\dd t = \frac1n u^{1/n - 1}\dd
u$:
\[
\int_0^1 \frac{\dd t}{\sqrt{1-t^n}}
= \frac1n\int_0^1 u^{\frac1n-1}(1-u)^{-1/2}\dd u
= \frac1n\,B\Bigl(\frac1n, \frac12\Bigr) .
\]
$n = 1$: $B\bigl(1,\frac12\bigr) = B\bigl(\frac12,1\bigr) = 2$,
coincidindo com $\int_0^1\frac{\dd t}{\sqrt{1-t}} = 2$. $n = 2$:
$\frac12 B\bigl(\frac12,\frac12\bigr) = \frac\pi2 = \arcsin 1$.
Para $n = 4$ o valor $\frac14 B\bigl(\frac14,\frac12\bigr)$ é
a constante da lemniscata: nenhuma forma fechada elementar.

\textbf{23.} Iterando a equação funcional:
\[
\frac{1}{\Gamma(x)}\int_0^\infty t^{x+k-1}\eu^{-t}\dd t
= \frac{\Gamma(x+k)}{\Gamma(x)}
= (x+k-1)(x+k-2)\cdots x ,
\]
o fatorial ascendente com $k$ fatores. Em $x = 1$:
$\Gamma(1+k)/\Gamma(1) = k!$, os momentos de $\eu^{-t}$ vindos do
\cref{exo:b2:integration:2}.

\textbf{24.} Substitua $t = \frac{1+s}2$ ($s \in
\intoo{-1}1$, $\dd t = \frac{\dd s}2$, $t(1-t) =
\frac{1-s^2}4$):
\[
B(x,x) = \int_{-1}^{1}\Bigl(\frac{1-s^2}{4}\Bigr)^{x-1}
\frac{\dd s}{2}
= 4^{1-x}\int_0^1 (1-s^2)^{x-1}\dd s
\]
(o integrando é par). Depois $s = \sqrt v$ ($\dd s =
\frac{\dd v}{2\sqrt v}$):
\[
B(x,x) = \frac{4^{1-x}}{2}\int_0^1
v^{-1/2}(1-v)^{x-1}\dd v
= 2^{1-2x}\,B\Bigl(\frac12, x\Bigr) .
\]
Para $2x \in \N^*$ todos os argumentos à vista estão em
$\frac12\N^*$, de modo que a fórmula de Euler (questão 14) se aplica aos dois
lados:
\[
\frac{\Gamma(x)^2}{\Gamma(2x)}
= 2^{1-2x}\,
\frac{\Gamma\bigl(\frac12\bigr)\Gamma(x)}
{\Gamma\bigl(x+\frac12\bigr)}
\quad\Longleftrightarrow\quad
\Gamma(x)\,\Gamma\Bigl(x+\frac12\Bigr)
= 2^{1-2x}\sqrt\pi\;\Gamma(2x) .
\]
Verificação direta em $x = n$: o membro da esquerda é $(n-1)!\cdot
\frac{(2n)!\sqrt\pi}{4^n n!} = \frac{(2n)!\sqrt\pi}{4^n n}$,
o da direita $2\cdot4^{-n}\sqrt\pi\,(2n-1)! =
\frac{(2n)!\sqrt\pi}{4^n n}$: iguais.

\textbf{25.} (i) A integração por partes produziu $B(x,y+1) =
\frac yx B(x+1,y)$, a única identidade da qual decorrem toda
relação de descida, os valores inteiros e semi-inteiros e a
recorrência de Wallis. (ii) A convergência dominada transformou
as integrais elementares $\int_0^n(1-t/n)^n t^{x-1}$ em
$\Gamma(x)$ (fórmula de Gauss, questão 17) e calculou o
limite $a \to \infty$ na questão 21. (iii) Do capítulo de
comparação importamos a constante de Euler ($\ln n - H_n \to
-\gamma$, questão 18) e a assintótica do binomial central
(questão 19) --- isto é, a fórmula de Stirling disfarçada. (iv)
A fórmula de Euler $B(x,y) = \Gamma(x)\Gamma(y)/\Gamma(x+y)$ está
agora demonstrada para $y \in \N^*$ com $x > 0$ arbitrário (questão
10) e para todos os pares semi-inteiros (questão 14); o caso
geral $x, y > 0$ aguarda o cálculo por integral dupla do
capítulo sobre integrais múltiplas, que fatoriza
$\Gamma(x)\Gamma(y)$ sobre um quarto de plano.
\end{solution}
